\chapter{보존·격리·확장·압축: 안정성–가소성 mechanism}
\label{bg:chapter-stability-methods}

이 장은 continual learning의 대표 도구를 ``무엇을 저장하고 어떤 update를 금지하는가''로 설명한다.
알고리즘 이름보다 중요한 것은 data, auxiliary state, constraint, capacity 회수 경로다.

\section{Regularization: 중요한 parameter는 덜 움직인다}

\concept{regularization}{정규화}는 새 task loss에 보존 조건을 더한다. EWC는 과거 parameter
$\theta_i^*$와 importance $F_i$를 저장하고 다음 objective를 사용한다
\parencite{kirkpatrick2017ewc}.

\begin{equation}
\mathcal{L}_{EWC}(\theta)=\mathcal{L}_{new}(\theta)+
\frac{\lambda}{2}\sum_i F_i(\theta_i-\theta_i^*)^2.
\end{equation}

\concept{ewc}{EWC}의 직관은 $F_i$가 큰 parameter를 ``과거에 중요하므로 비싸게 움직이는
spring''으로 만드는 것이다. raw replay가 없어도 되지만 매 cycle마다 reference weights와 importance를
보존해야 한다. task가 늘수록 importance를 합칠지, 최근 것만 둘지, 충돌하는 중요도를 어떻게 낮출지
결정해야 한다. LLM 전체 weight에 per-parameter FP32 importance를 두면 state가 model size만큼 늘 수
있으므로 layer/block/low-rank 단위 근사가 필요하다.

\section{Gradient constraint: 과거 손실을 늘리는 방향을 피한다}

새 gradient $g$와 replay example의 과거 gradient $g_k$가 있을 때 $g^\top g_k<0$이면 두 objective가
국소적으로 충돌한다. \concept{gradient-projection}{기울기 투영}은 update를

\begin{equation}
\tilde g=\arg\min_z \frac12\lVert z-g\rVert_2^2
\quad \text{s.t.}\quad z^\top g_k\ge 0\quad \forall k
\end{equation}

같은 feasible region으로 옮긴다. \concept{gem}{GEM}은 episodic memory와 이 제약을 결합해
negative backward transfer를 줄이고 positive transfer도 측정했다 \parencite{lopezpaz2017gem}.
그러나 reference gradient를 많이 둘수록 backward와 projection solve 비용이 커진다. Sleep cluster에서
compute를 쓸 수 있어도, 사용자 수에 비례한 gradient memory를 HBM에 올리는 것이 병목일 수 있다.

\section{Isolation과 growth: 충돌할 parameter를 나눈다}

\concept{parameter-isolation}{파라미터 격리}는 base를 freeze하고 user/task adapter, mask, expert를
따로 둔다. locality와 rollback이 쉽고 global forgetting을 줄이지만, active artifact 수가 늘며 routing
metadata와 random read가 증가한다. 같은 pattern을 여러 adapter가 중복 학습하면 logical capacity보다
physical capacity가 빠르게 소진된다.

\concept{dynamic-expansion}{동적 확장}은 새 분포가 기존 subspace로 충분히 표현되지 않을 때 module을
추가한다. 이는 새 data를 강제로 기존 weights에 넣는 것보다 plastic하지만, append-only라면 결국
capacity를 초과한다. expansion 조건에는 validation gain뿐 아니라 expected reuse와 bytes·latency가
들어가야 한다.

반대 방향의 \concept{pruning}{프루닝}은 importance가 낮은 weight, expert, memory item을 제거한다.
Parameter magnitude만 보면 희귀하지만 critical한 기억을 지울 수 있다. Sleep system의 prune unit은
``작은 수치''보다 provenance cluster, last verified time, future query utility, replacement coverage를
포함해야 한다.

\section{Capacity는 bytes만이 아니라 interference와 운영 한도다}

\concept{capacity-budget}{용량 예산} $B$를 단순 저장 bytes가 아니라 다음 resource vector로 둔다.

\begin{equation}
\mathbf{B}=(B_{bytes}, B_{latency}, B_{energy}, B_{risk}, B_{ops}).
\end{equation}

$B_{ops}$는 version 수, validation 시간, rollback 복잡도 같은 운영 capacity다. 외부 vector store가
남은 SSD bytes를 충분히 가져도 ANN latency와 index rebuild가 먼저 한도에 닿을 수 있다. LoRA file이
작아도 사용자 수가 수억이면 adapter lookup과 batching fragmentation이 한도다.

\concept{compression}{압축}은 허용 왜곡 안에서 representation cost를 줄인다. 여러 episode를 summary로,
여러 vector를 prototype으로, 여러 adapter를 merged adapter로 바꿀 수 있다. \concept{rate-distortion}{율–왜곡}
관점에서는 memory representation $Z$의 rate $R(Z)$와 future query distribution에서의 distortion
$D(X,\hat X)$를 함께 최적화한다.

\begin{equation}
\min_{q(z|x)}\; \mathbb{E}[D(X,\hat X(Z))]+\beta R(Z).
\end{equation}

여기서 $D$는 원문 reconstruction loss일 필요가 없다. future task success, factual calibration,
provenance recovery, deletion granularity를 함께 반영해야 한다. 이 framework를 memory compaction에
직접 연결하는 연구가 등장하고 있으나 \parencite{ratedistortion2026}, 어떤 distortion이 장기 agent
utility를 대변하는지는 아직 열린 측정 문제다.

\section{Admission, eviction, consolidation은 하나의 정책이다}

\concept{memory-admission}{기억 입장 정책}은 새 event를 discard, raw log, text summary, vector,
graph, adapter candidate, global training queue 중 어디로 보낼지 정한다. \concept{memory-eviction}{기억
퇴거 정책}은 예산이 부족할 때 delete, compress, demote, merge 중 하나를 고른다. 둘을 따로 최적화하면
새 item을 받아 놓고 즉시 비슷한 old item을 지우는 thrashing이 생긴다.

\concept{consolidation}{공고화}는 여러 경험을 더 안정적인 표현으로 결합하는 전체 변환이다. 중요한
점은 destination이 항상 model weights가 아니라는 것이다. episode $\rightarrow$ temporal graph,
vector cluster $\rightarrow$ prototype, repeated correction $\rightarrow$ LoRA candidate처럼 evidence와
reuse에 따라 경로가 달라진다.

\begin{table}[htbp]
\centering
\caption{Bounded memory에서 쓰는 네 가지 capacity 동작}
\label{tab:bg-capacity-actions}
\small
\begin{tabularx}{\textwidth}{p{0.14\textwidth} X X X}
\toprule
동작 & 얻는 것 & 잃을 수 있는 것 & 필수 검증 \\
\midrule
delete & 즉시 bytes·rights 회수 & future utility, tail evidence & 삭제 대상·파생물 완전성 \\
compress & lower tier 비용 감소 & 세부·불확실성·provenance & distortion slice별 회귀 \\
merge/distill & 중복 compute 감소 & 충돌 사실, minority behavior & locality와 conflict test \\
promote & wake latency 감소 가능 & weight interference, rollback 비용 & retention·safety·break-even \\
\bottomrule
\end{tabularx}
\end{table}

\begin{keypoint}[단조 증가를 막는 invariant]
각 sleep cycle은 새 bytes만 보고하지 말고 $\Delta B$를 destination별로 보고해야 한다. admission된
상태가 있으면 eviction·compression·reuse gain 중 무엇이 그 비용을 정당화하는지 연결한다. ``모두
보관 후 나중에 정리''는 retention policy가 아니라 unbounded debt다.
\end{keypoint}
